% chapters/ch15_project_deepdives.tex — Chapter 15: Her Project Deep-Dives
% GATED CHAPTER. Per PLAN.md Section 5 and CLAUDE.md, this chapter is written
% ONLY from inputs/INTAKE.md. Until its REQUIRED fields are filled, this file is
% a SKELETON with [NEEDS INTAKE: ...] placeholders and contains NO fabricated
% project facts (no invented defects, metrics, tools, roles, or numbers).
\chapter{Her Project Deep-Dives}
\label{ch:project-deepdives}

\begin{partnote}
\textbf{GATED --- skeleton only.} Per \code{PLAN.md} Section 5, every project
fact in this chapter (what was tested, tools used, defects found, metrics, team
size, exact role) must come from \code{inputs/INTAKE.md} --- never invented. The
REQUIRED fields of \code{inputs/INTAKE.md} are not yet filled, so below is the
\emph{structure} only, with \textbf{[NEEDS INTAKE: ...]} markers showing exactly
what to supply. \textbf{Do not improvise these in an interview from the
skeleton} --- a fabricated detail repeated confidently is dismantled by one
follow-up. Fill INTAKE.md, then this chapter is generated in full: for each
project, a 2-minute story, a 5-minute deep version, a block diagram to draw,
15 likely follow-ups with answers, and prepared \emph{honest-boundary} lines.
\end{partnote}

For each project the generated chapter will contain: \textbf{(a)} a 2-minute
spoken story (Situation/role $\rightarrow$ what you did $\rightarrow$ result);
\textbf{(b)} a 5-minute deep version; \textbf{(c)} a block diagram you can draw
on a whiteboard; \textbf{(d)} 15 likely follow-up questions \emph{with} your
answers, including the uncomfortable ones (why this approach? what failed? exact
latency/baud/voltage numbers?); and \textbf{(e)} 2--3 prepared honest-boundary
sentences for questions beyond your scope. The technical framing for each is
already in the chapters cross-referenced below; only the \emph{personal facts}
are missing.

\begin{keyidea}
\textbf{Keep the two employers' domains separate} --- this is what makes every
answer consistent. \textbf{TASL (current, UAV R\&D):} UAV DFCC \& signal-
conditioning cards, UAV LRU integration, board bring-up, RS422 Embedded C
interface development, multi-protocol bring-up. \textbf{BEL (LCA Tejas,
flight-critical):} DFCC verification, CCDL, SPIL, DPRAM, MIL-STD-1553B Python
automation, DO-178B Level A. Never let CCDL/SPIL drift into the TASL story or
the UAV board bring-up into the BEL story.
\end{keyidea}

%=====================================================================
\section{TASL --- UAV R\&D (Current Role)}
%=====================================================================

\subsection{Project 1 --- UAV DFCC Card Validation}
\emph{Framing: Chapters 8 (avionics buses), 11 (bring-up/debug), 12 (V\&V).}
\begin{itemize}
\item \textbf{[NEEDS INTAKE: Q1 --- exactly what signals/parameters do you check
  on a UAV DFCC card, and the pass/fail criteria?]}
\item \textbf{[NEEDS INTAKE: Q2 --- one real defect you found and how it was
  resolved (symptom, root cause, fix, how verified)?]}
\end{itemize}

\subsection{Project 2 --- AETS-DFCC Signal-Conditioning Card Integration}
\emph{Framing: Chapters 11 (lab/debug), 10 (signal/sensor), 8.}
\begin{itemize}
\item \textbf{[NEEDS INTAKE: Q3 --- which sensor/analog signals does the SC card
  condition, and what do you verify on the acquisition/conditioning path?]}
\item \textbf{[NEEDS INTAKE: Q4 --- which instruments (oscilloscope, DMM,
  electronic load) and what measurements?]}
\end{itemize}

\subsection{Project 3 --- UAV LRU Integration \& Board Bring-up}
\emph{Framing: Chapters 11 (bring-up sequence), 3 (boards/MCU), 4--8 (interfaces).}
\begin{itemize}
\item \textbf{[NEEDS INTAKE: Q5 --- which boards do you bring up (autopilot
  v2.2/4.0/4.1/6.1, PDU, IRNSS, VMNSC, baseboard) and what's your bring-up
  sequence with ST-Link/J-Link?]}
\item \textbf{[NEEDS INTAKE: Q6 --- which LRUs do you integrate (PSU200, AU, AIU,
  INSU, CAN servos, radio, pitot/Laversab) and what does integration involve?]}
\end{itemize}

\subsection{Project 4 --- RS422 Embedded C Interface Development}
\emph{Framing: Chapters 1 (Embedded C), 5 (RS422).}
\begin{itemize}
\item \textbf{[NEEDS INTAKE: Q7 --- what do the RS422 Embedded C interface/test
  routines you write actually do (interface code? test harness? what function)?]}
\item \textbf{[NEEDS INTAKE: Q8 --- roughly how much code, on what target, and
  how is it verified?]}
\end{itemize}

%=====================================================================
\section{BEL --- LCA Tejas, Flight-Critical Avionics (Previous Role)}
%=====================================================================

\subsection{Project 5 --- CCDL Verification}
\emph{Framing: Chapter 8 (CCDL, voting, failover), Chapter 10 (redundant control).}
\begin{itemize}
\item \textbf{[NEEDS INTAKE: Q9 --- what specifically do you verify on the CCDL
  (which data, what timing bound, what failover/redundancy behaviour), and how?]}
\item \textbf{[NEEDS INTAKE: Q10 --- any concrete number: frame rate, latency
  bound, channel count? Only if real.]}
\end{itemize}

\subsection{Project 6 --- SPIL Verification}
\emph{Framing: Chapters 5 (serial/RS422), 8 (serial inter-card links).}
\begin{itemize}
\item \textbf{[NEEDS INTAKE: Q11 --- what does the SPIL link carry, and what do
  you test at the packet level (framing, integrity, timing, loss behaviour)?]}
\item \textbf{[NEEDS INTAKE: Q12 --- the tools and the exact diagnostics you run?]}
\end{itemize}

\subsection{Project 7 --- DPRAM Architecture Validation}
\emph{Framing: Chapter 9 (DPRAM handshake), 4 (shared data), 3 (ARM9/LPC3250).}
\begin{itemize}
\item \textbf{[NEEDS INTAKE: Q13 --- the real-time DPRAM polling architecture you
  validated: what data/command handshakes, what made it sub-1\,ms, your role?]}
\item \textbf{[NEEDS INTAKE: Q14 --- the white-box firmware testing on LPC3250
  (ARM9): what you actually tested and found; the BSP/RTSM/driver work?]}
\end{itemize}

\subsection{Project 8 --- MIL-STD-1553B Python Automation}
\emph{Framing: Chapters 14 (Python automation), 8 (1553B), 12 (V\&V).}
\begin{itemize}
\item \textbf{[NEEDS INTAKE: Q15 --- the Python suite: what it parsed (1553B /
  RS422), how it was built, how the about-60\% number was measured (baseline vs
  after)?]}
\end{itemize}

%=====================================================================
\section{Claim--Resume Consistency Check}
%=====================================================================
\begin{partnote}
When INTAKE.md is filled, every claim written here must be cross-checked against
the resume source (the \code{.tex} in \code{resume/}) per PLAN.md Section 5.4. If
a resume bullet overstates what INTAKE.md reveals, flag it in \code{progress.md}
so the \emph{resume} is corrected --- the resume bends to reality, never the
reverse. No such check can be completed until the REQUIRED INTAKE fields exist;
this is recorded as an open gate in \code{progress.md}.
\end{partnote}
